\subsection{Implications of Our Results and Prior Work}
\label{sec:results}		

Below we describe in further detail the implications of our results, and discuss previous work on MOL problems. In the following problems, general matroids are assumed to be accessed via membership oracles. 

\paragraph{Exact Matroid Basis (EMB):}
This is a generalization of the $k$-{\sc subset sum} problem (where $(E,\cI)$ is a uniform matroid).\footnote{In a uniform matroid, $\cI= \{S \subseteq E~|~ |S| \leq k\}$.} Thus, EMB is unlikely to be solvable in polynomial time. Instead, we seek a pseudo-polynomial time algorithm whose running time has polynomial dependence on the encoding size of the instance and the target value $T$. 
Indeed, the special case of EMB in which the matroid is {\em representable} (or, linear) admits such  a pseudo-polynomial time algorithm \cite{CGM92}. Since the 1990s, it has been an open question whether the result of Camerini et al.~\cite{CGM92} can be extended to general matroids. 
\Cref{thm:emb} resolves this question, ruling out the existence of  a pseudo-polynomial time algorithm for EMB.
%Using our hardness results for EMB, we can potentially show that a parameterized matroid problem, studied by Fomin et al.~\cite{FGKSS23} on representable matroids, cannot be solved efficiently on general matroids (we leave it for a future work, see \Cref{sec:discussion}). 

\paragraph{Budgeted Matroid Independent Set (BM):} 

This problem is cast as $(\max,\IS,\leq)$-MOL. BM is a natural generalization of the classic $0/1${\sc-knapsack} problem, for which a Fully PTAS has been known since the 1970s~\cite{La79}. As mentioned above, a Fully PTAS is known also for other special cases of BM. A PTAS for BM was first given in \cite{BBGS11} as a special case of {\sc budgeted matroid intersection (BMI)}. In this generalization of BM, we are given {\em two} matroids $(E,\cI_1)$ and $(E,\cI_2)$, and a solution has to be an independent set of {\em both} matroids. A PTAS for BM also follows from 
the results of~\cite{GZ10,CVZ11} which present PTASs for multi-budgeted variants of BM. An Efficient PTAS for BM was recently given in \cite{DKS23} and for BMI in \cite{DKS23b}. The existence of a Fully PTAS for BM was posed as a central open question in~\cite{BBGS11,DKS23,DKS23b}. We answer this question negatively, as formalized in \Cref{thm:main}, giving a tight lower bound for BM and BMI. %Our results can be used to obtain a lower bound also for a {\em parameterized} version of BM \cite{doron2023budgeted} (see \Cref{sec:discussion}). 

%A slight modification\Cref{thm:bm} rules out a polynomial {\em parametrized approximation scheme} with running time polynomial in $\frac{1}{\eps}$ for a parametrized variant of BM studied in~\cite{doron2023budgeted}, where $\eps$ is the error parameter.  %the next result.
%\begin{restatable}{theorem}{thmbm}
%	\label{thm:bm}
%	There is no randomized \textnormal{Fully PTAS}
%	%fully polynomial time approximation scheme \textnormal{(FPTAS)}  
%	for {\sc budgeted matroid independent set}.
%\end{restatable} 

%\Cref{thm:bm} does not use any complexity related assumptions. The proof is information theoretic and exploits the limitations of the oracle model. As mentioned
%above, we also prove a variant of \Cref{thm:bm} which rules out the existence of a Fully PTAS for a version of BM in which the matroid belongs to a specific family of matroids and is encoded as part of the input for the algorithm. 


%	 	Observe that if  the requirement that the solution is an independent set is removed from BM  (equivalently, $\cI=2^E$), the problem becomes the classic 
%	 	$0/1$-knapsack, which admits a Fully PTAS. Furthermore, if the budget constraint is removed (for example, by setting $c(e)=0$ for all $e\in E$) then we have the  maximum weight matroid independent set problem, which is solvable in polynomial time. That is, the bound in \Cref{thm:bm} is stronger than any hardness of approximation bound which can be derived for one of these special cases, and hence must leverage the interaction between the matroid  and the budget constraints. The same holds for the other
%	 	problems considered in this paper. This distinguishes our results from
%	 	 existing lower bounds for matroid problems (e.g.,~\cite{JK82}).


%\Cref{thm:bm} rules out the existence of a Fully PTAS also for the following well studied generalization of BM.

\comment{

\problem{{\sc Budgeted Matroid Intersection} (BMI)}{
	{\bf Instance} &  
	$I = (E,\cI_1, \cI_2,c,p,B)$, where $(E,\cI_1)$ and $(E,\cI_2)$ are  matroids,
	$c:E \rightarrow \mathbb{N}$ is a cost function, $p:E \rightarrow \mathbb{N}$ is a profit function and $B \in \mathbb{N}$ is the budget.
	\\
	{\bf Solution}& $S\in \cI_1\cap \cI_2$ such that $c(S) \leq B$.\\
	{\bf Objective} & Find a solution $S$ for $I$ of maximum profit $p(S)$. \\
	{\bf Model } &Both $\cI_1$ and $\cI_2$ are accessible only via membership oracles; the encoding size of the instance $I$ is the encoding size of $(E,c,p,B)$, the instance without $\cI_1$ and~$\cI_2$.}
}

%That is, BMI is a variant of BM in which we are given {\em two} matroids $(E,\cI_1)$ and $(E,\cI_2)$, and a solution has to be an independent set of {\em both} matroids. 
%A PTAS for BMI was first given in~\cite{BBGS11}, and recently improved to an Efficient PTAS \cite{DKS23b}. Since BMI is a generalization of BM (e.g., take $\cI_1=\cI_2$), \Cref{thm:bm} implies the following.
%\begin{cor}
%	\label{cor:bmi}
%	There is no randomized  \textnormal{Fully PTAS}
%	%fully polynomial time approximation scheme \textnormal{(FPTAS)}  
%	for {\sc budgeted matroid intersection}.
%	% \textnormal {(BMI)}.
%\end{cor}
%That is, the Efficient PTAS of \cite{DKS23b} for BMI is the best possible.

%\subsubsection{Exact Weight Matroid Independent Set}

\comment{
As  in \Cref{thm:bm},  \Cref{thm:emi} does not require any complexity related assumptions. In particular, \Cref{thm:emi} implies that there is no algorithm for EMI which runs in time $(T\cdot |I|)^{O(1)}$.

%A pseudo polynomial for a variant of EMI in which the matroid is {\em representable}  follows from \cite{CGM92} and \cite{LMPS18}. 
Camerini et al.~\cite{CGM92} studied EMI on {\em representable} matroids. A {\em representable matroid} over a field $\mathbb{F}$ is a matroid $(E,\cI)$, where $E\subseteq \mathbb{F}^{\ell}$ for some ${\ell}\in \mathbb{N}$, and $$\cI=\{S\subseteq E~|~\textnormal{$S$ is linearly independent in $\mathbb{F}^{\ell}$}\}.$$
While it is easy to show the above construction always leads to a matroid, not all matroids are representable (see, e.g.,~\cite{N16}).  
%\cite{Ox2006}) 
%We replaced by a ref to Nelson's paper. --H
%\ariel{verify this claim indeed appears there}.
Assuming the elements of $\mathbb{F}$ can be encoded,
representable matroids can be easily encoded within a problem instance by simply encoding the vectors in $E$ as part of the input. 

 The variant of EMI studied by \cite{CGM92} considers representable matroids over the field of rational numbers. Specifically, the input for {\em representable} EMI is a representable matroid  $(E,\cI)$ over the rationals, a weight function $w:E\rightarrow \mathbb{N}$, and a target weight $T\in \mathbb{N}$. The objective is to determine if there exists an independent set  of the matroid of weight exactly $T$, that is, $S\in \cI$ such that $w(S)=T$. The result of \cite{CGM92} together with \cite{LMPS18} yield a randomized algorithm for representable EMI whose running time is polynomial
in the encoding of the instance and in $T$.\footnote{The result in \cite{CGM92} requires the solution to be a basis and is given for the more generic {\em exact parity basis} problem. The requirement that the solution is a basis can be easily circumvented using the truncation technique of \cite{LMPS18} for representable matroids.} \Cref{thm:emi} implies that this result cannot be generalized to arbitrary matroids. 

\Cref{thm:emi} can also be used to show that a parameterized matroid problem, studied by Fomin et al.~\cite{FGKSS23} on representable matroids, cannot be solved efficiently on general matroids. This additional application uses the fact that \Cref{thm:emi} also rules out algorithms
whose running times depend on a parameter $k$. 

}

%\Cref{thm:emi} can also be used to show that an algorithm of Fomin et al.\ \cite{FGKSS23}, which solves a parameterized matroid problem for representable matroids, cannot be generalized to arbitrary matroids. We give the details in \Cref{sec:additional}. This additional application uses the fact that \Cref{thm:emi} also rules out algorithms
%whose running times depend on a parameter $k$. 

\paragraph{Constrained Minimum Basis of a Matroid (CMB):} This problem can be cast as the matroid optimization problem $(\min,\bases,\leq)$-MOL. 
The {\sc constrained minimum spanning tree (CST)} problem is the  special case of CMB in which the matroid $(E,\cI)$ is graphical \cite{RG96, andersen1996bicriterion,HL04,hong2004fully}, namely, there is a graph $G=(V,E)$ such that  
the independent sets $\cI$ are cycle-free subsets of edges in $G$.
A PTAS for CST was given by Ravi and Goemans~\cite{RG96}. This result was improved to an Efficient PTAS by Hassin and Levin \cite{HL04}. A {\em bicriteria} FPTAS, which violates the budget constraint by a factor of $(1+\eps)$, was presented in~\cite{hong2004fully}. The authors of~\cite{HL04} mention that their result actually gives an Efficient PTAS for CMB. 
The existence of a Fully PTAS for CMB remained an open question. \Cref{thm:main} shows that the Efficient PTAS for CMB cannot be improved.
%This fits better for the Discussion. --H 
%It remains an open question whether CST admits a Fully PTAS.
 
\paragraph{Knapsack Cover with a Matroid (KCM):}
As a final implication, \Cref{thm:main} rules out the existence of a Fully PTAS for a  {\em coverage} variant of $0/1${\sc-knapsack}, formulated as $(\min,\IS,\geq)$-MOL. 
In \cite{CCNR13}, Chakaravarthy et al. presented a PTAS for KCM using integrality properties of a linear programming formulation of KCM. Moreover, for the special case of KCM with a {\em partition matroid}, they give a Fully PTAS based on dynamic programming. The existence of  a Fully PTAS for KCM  on a general matroid was posed in~\cite{CCNR13} as an open question. \Cref{thm:main} answers this question negatively.
Our initial study indicates that an Efficient PTAS  for KCM can potentially be obtained by adapting the approach of Hassin and Levin \cite{HL04} to the setting of KCM. This suggests that our lower bound cannot be strengthened. 


\paragraph{Organization:}
\label{sec:organization}

In \Cref{sec:technical} we introduce the $\Pi$-matroid family and give the proof of \Cref{thm:emb}. In \Cref{sec:tight_absolute} we prove \Cref{thm:main}, and in
 \Cref{sec:SAT} we show that similar lower bounds hold in the standard computational model. We conclude in \Cref{sec:discussion} with a summary and directions for future work.